\section{问题分析}

\subsection{问题一的分析}

题目已给出交会定位法：每个检测点的示向度存在左右各 $1^\circ$ 的误差范围，多个检测点对应区域的公共部分即为干扰源的定位区域。因此，本问首先需要从几何上说明，在有效交会且公共区域有界的情况下，多个区域交汇所得的定位区域必为凸多边形。

定位区域的直径定义为区域内任意两点之间距离的最大值。对于凸多边形，需要进一步证明这一最大值必能在两个顶点之间取得，从而将区域直径转化为最远顶点对之间的距离。得到直径端点后即可作圆，但该圆不一定覆盖整个定位区域，还需建立相应的几何判据。利用定位区域与圆的凸性，只需判断凸多边形的全部顶点是否位于圆内，即可确定该圆能否覆盖整个定位区域。

\subsection{问题二的分析}

问题一由交会定位法得到多边形定位区域。然而在实际情形下，定位区域可能由于目标区域有界及信号存在有效接收半径，从而产生圆弧边。因此，问题二首先对含圆弧的可行域进行保守多边形化：以圆域的外切正多边形构造粗外包，使后续计算仍能沿用凸多边形方法，并保证不漏掉真实源。

问题一还说明，以区域直径为直径作圆并不一定覆盖整个定位区域，故单独使用直径的一半不能作为可靠的定位误差上界。为此，问题二统一用定位区域的最小包围圆半径定义“定位不确定性半径”。选点时，在区域内取若干有代表性的目标位置，分别预测第二次定位后的不确定性半径，再取其中最大的一个，作为主要效果指标。

第二测点的选择需要同时考虑接收是否可靠、两次测向形成的交会角度是否理想，以及机器狗到达该点的行动时间。本文先取一批离散候选点做综合打分，再在保证接收区域内，以 Fisher 信息为准则做连续选点，并限制机器狗的行动时间。连续方法产生交会几何质量较好的候选点，所有候选点再按相同的定位不确定性半径指标复评分，并给出按行动时间和定位不确定性半径定义的非支配候选集，再按分级比较规则确定推荐测点；同时将保证能接收到信号的区域作为候选区域。

